\scanpage{179}
% archive foliation: 178
D'autre part $\beta$ est un mono. sous \unsure{les hypothèses|l'hypothèse} $\underline H^0_{X/S}(\underline O_X)=\underline O_S$ ``universelle''; et il est surjectif donc bijectif si on suppose \uline{de plus} que \deleted{$\unsure{\operatorname{Pic}}$} $X$ a localement une section sur $S$.

\uline{Exemple} où $\beta$ n'est pas surjectif: variété de \unsure{Séveri}-Brauer non triviale sur un corps .... On verra qu'il y a une obstruction à la surjectivité de $\beta$, à valeurs dans un $H^2(S,\mathbf G_m)$ .....

\uline{Remarque}

On voit donc que le foncteur \uuline{Prépic} n'est pas représentable \deleted{en général, \unsure{que dans un}} \inserted{sauf dans des cas} \deleted{\unsure{peu}} triviaux sans intérêt; $\mathrm{Prépic}'_{X/S}$ ne peut être représentable que si $\beta$ est un isom et si $\underline{\underline{\operatorname{Pic}}}_{X/S}$ est représentable, \deleted{\unsure{(que si} $\unsure{\beta}$} [ce qui est pratiquement lié à des conditions comme celles déjà envisagées, à savoir l'existence d'une section de $X$ sur $S$]. Quant à $\underline{\underline{\operatorname{Pic}}}_{X/S}$, il est représentable dans des cas généraux importants qui seront étudiés par la suite. C'est \unsure{peut être} le plus \unsure{important pour notre propos}.
